\documentclass[11pt,a4paper]{article}
\usepackage[utf8]{inputenc}
\usepackage[T1]{fontenc}
\usepackage{amsmath,amssymb,amsthm}
\usepackage{physics}
\usepackage{geometry}
\usepackage{booktabs}
\usepackage{array}
\usepackage{hyperref}
\usepackage{mathtools}
\usepackage{setspace}
\usepackage{enumitem}
\usepackage{fancyhdr}
\usepackage{tcolorbox}
\tcbuselibrary{theorems,skins,breakable}

\geometry{margin=2.6cm, top=3.2cm, bottom=3.2cm}
\onehalfspacing

\hypersetup{
  colorlinks=true,
  linkcolor=black,
  citecolor=black,
  urlcolor=black,
  pdftitle={QGD Chapter 4: Equations of Motion},
  pdfauthor={Romeo Matshaba}
}

\pagestyle{fancy}
\fancyhf{}
\fancyhead[L]{\small\textit{Quantum Gravitational Dynamics}}
\fancyhead[R]{\small\textit{Chapter 4: Equations of Motion}}
\fancyfoot[C]{\thepage}
\renewcommand{\headrulewidth}{0.4pt}

\theoremstyle{plain}
\newtheorem{theorem}{Theorem}[section]
\newtheorem{lemma}[theorem]{Lemma}
\newtheorem{corollary}[theorem]{Corollary}
\newtheorem{proposition}[theorem]{Proposition}
\theoremstyle{definition}
\newtheorem{definition}[theorem]{Definition}
\theoremstyle{remark}
\newtheorem{remark}[theorem]{Remark}

\newcommand{\sigt}{\sigma_t}
\newcommand{\lQ}{\ell_Q}
\newcommand{\mQ}{m_Q}
\newcommand{\boxg}{\square_g}
\newcommand{\half}{\tfrac{1}{2}}
\newcommand{\Mpl}{M_{\mathrm{Pl}}}

\title{%
  \vspace{-1.2cm}
  {\large\scshape Quantum Gravitational Dynamics}\\[6pt]
  {\LARGE\bfseries Chapter 4: Equations of Motion}\\[8pt]
  {\normalsize From the $\sigma$-Field to the Worldline}
}
\author{%
  Romeo Matshaba\\[4pt]
  \normalsize Department of Physics, University of South Africa
}
\date{March 2026}

\begin{document}
\maketitle
\thispagestyle{empty}

\begin{abstract}
\noindent
We derive the equations of motion in Quantum Gravitational Dynamics at three levels
of description: (i) the $\sigma$-field equation governing the entire system; (ii) the
modified worldline equation for a test particle, which replaces the geodesic equation
of General Relativity; and (iii) the conserved stress-energy equation encoding the
interplay between matter, the $\sigma$-field, and the Pais--Uhlenbeck massive mode.
The central structural result is the Advection Tensor $\Omega^\lambda_{\mu\nu}$, which
arises from the composite nature of the QGD metric $g_{\mu\nu}(\sigma_t)$ and encodes
the field-dragging effect absent in GR.  We derive $\Omega^\lambda_{\mu\nu}$ exactly
from the particle action, compute its explicit components in the Schwarzschild sector,
and establish the precise sense in which it generalises and reduces to the standard
Levi-Civita connection.  The relationship between the three levels is made explicit,
and a comparison table between GR and QGD motions is provided.
\end{abstract}

\tableofcontents
\newpage

%=============================================================
\section{Overview: Three Levels of Motion in QGD}
\label{sec:overview}
%=============================================================

In General Relativity the motion of matter is governed by two logically separate
laws: the geodesic equation (worldline) and the Einstein field equations (geometry).
In QGD these are unified.  The metric is a composite functional of the $\sigma$-field,
$g_{\mu\nu} = g_{\mu\nu}(\sigma_t)$, so the field equation and the worldline equation
are both consequences of the same master action.

The equations of motion appear at three levels, each a projection of the same
underlying field theory:

\begin{enumerate}
\item \textbf{Field level.}  The Pais--Uhlenbeck master equation governs the
  $\sigma$-field sourced by all matter:
  \begin{equation}
    \boxg\!\left(1 - \kappa\lQ^2\,\boxg\right)\sigma_\mu = S_\mu,
    \qquad S_\mu = Q_\mu + G_\mu + T_\mu.
    \label{eq:PU_field}
  \end{equation}
  This is the most fundamental statement of motion in QGD.

\item \textbf{Worldline level.}  For a test particle of mass $\mu$, varying the
  particle action $S_{\rm pp} = -\mu\int ds$ with respect to the worldline
  $x^\mu(\lambda)$ in the composite metric $g_{\mu\nu}(\sigma_t)$ gives:
  \begin{equation}
    \frac{d^2 x^\mu}{ds^2}
    + \left(\Gamma^\mu_{\alpha\beta} + \Omega^\mu_{\alpha\beta}\right)
      \frac{dx^\alpha}{ds}\frac{dx^\beta}{ds} = 0,
    \label{eq:modified_geodesic}
  \end{equation}
  where $\Gamma^\mu_{\alpha\beta}$ is the standard Levi-Civita connection of the
  background field $g_{\mu\nu}(\sigma_2)$ and $\Omega^\mu_{\alpha\beta}$ is the
  \emph{Advection Tensor} arising from the composite metric structure.

\item \textbf{Conservation level.}  The stress-energy of the full system---matter,
  $\sigma$-field, and Pais--Uhlenbeck massive mode---is covariantly conserved:
  \begin{equation}
    \nabla_\nu\!\left(
      T^{\mu\nu}_{\rm matter}
      + T^{\mu\nu}_{\rm field}
      + \Theta^{\mu\nu}_{\rm PU}
    \right) = 0.
    \label{eq:full_conservation}
  \end{equation}
\end{enumerate}

In the limit $\lQ \to 0$ and for a static, weak $\sigma$-field, all three levels
reduce to the standard GR results.  The QGD corrections---the Advection Tensor at
the worldline level and the PU stress tensor at the conservation level---become
non-negligible only at post-Newtonian order or near the Planck scale.

%=============================================================
\section{The Field-Level Equation of Motion}
\label{sec:field_eom}
%=============================================================

\subsection{Particles as field configurations}

In QGD a particle of mass $m_a$ is not a structureless point on a background
geometry; it is identified with its own single-body $\sigma$-field configuration,
\begin{equation}
  \sigt^{(a)}(\mathbf{x}) = \sqrt{\frac{r_s^{(a)}}{|\mathbf{x} - \mathbf{x}_a|}},
  \qquad r_s^{(a)} \equiv \frac{2Gm_a}{c^2}.
  \label{eq:single_body_field}
\end{equation}
This is the exact Schwarzschild solution for body $a$ in isolation.  The ``position''
of the particle is the center of mass of this field configuration, and its ``motion''
is the evolution of that center of mass as the field adjusts to the sources present.

The total field for an $N$-body system is the superposition plus interactions:
\begin{equation}
  \sigt^{\rm tot}(\mathbf{x},t)
  = \sum_{a=1}^N \sigt^{(a)}(\mathbf{x}-\mathbf{x}_a(t))
    + \delta\sigma(\mathbf{x},\{\mathbf{x}_a\},\{\mathbf{v}_a\}).
\end{equation}
The interaction field $\delta\sigma$ is determined by the integral equation from
Chapter~3: $\sigma = \sigma^{\rm free} + \int[G_0 - G_{\mQ}]S$.

\subsection{The self-consistent field equation}

The motion of body $a$ is encoded in the requirement that the master field
equation~\eqref{eq:PU_field} is satisfied at all times.  Projected onto the worldline
of body $a$, this becomes a self-consistency condition: the field $\sigt^{(a)}$ centred
on $\mathbf{x}_a(t)$ must satisfy the field equation in the background of all other
bodies.  There is no separate ``geodesic equation'' at the fundamental level.

The trajectory $\mathbf{x}_a(t)$ is determined implicitly by
\begin{equation}
  \boxg\!\left(1 - \kappa\lQ^2\,\boxg\right)
  \sigt^{(a)}\bigl(\mathbf{x} - \mathbf{x}_a(t)\bigr)
  = -Q_\mu^{\rm cross}\!\left[\sigt^{(a)}, \sum_{b\neq a}\sigt^{(b)}\right] + \cdots,
  \label{eq:self_consistent}
\end{equation}
where $Q_\mu^{\rm cross}$ is the cross-source from the nonlinear self-interaction
(Section~4.1 of Chapter~3).

%=============================================================
\section{The Modified Geodesic Equation}
\label{sec:modified_geodesic}
%=============================================================

\subsection{Derivation from the particle action}

For a test particle (a body whose own $\sigma$-field does not appreciably back-react
on the background), the effective worldline is derived by varying the particle action
\begin{equation}
  S_{\rm pp} = -\mu\int \sqrt{-g_{\mu\nu}(\sigt)\,\dot x^\mu\,\dot x^\nu}\;d\lambda
  \label{eq:particle_action}
\end{equation}
with respect to $x^\mu(\lambda)$.  The Euler--Lagrange equations give the standard-looking
form with modified Christoffel symbols.

The crucial step is the evaluation of $\partial g_{\alpha\beta}/\partial x^\mu$.  Since
$g_{\alpha\beta} = g_{\alpha\beta}(\sigt)$ and $\sigt = \sigt(\mathbf{x})$ is a field, the
chain rule gives:
\begin{equation}
  \frac{\partial g_{\alpha\beta}}{\partial x^\mu}
  = \frac{\partial g_{\alpha\beta}}{\partial\sigt}\,\nabla_\mu\sigt.
  \label{eq:metric_chain}
\end{equation}
Substituting~\eqref{eq:metric_chain} into the definition of the Christoffel symbol
$\Gamma^\lambda_{\mu\nu} = \half g^{\lambda\rho}(\partial_\mu g_{\rho\nu}
+ \partial_\nu g_{\rho\mu} - \partial_\rho g_{\mu\nu})$ yields the exact QGD
connection:
\begin{equation}
  \boxed{
  \Gamma^\lambda_{\mu\nu}[\sigt]
  = \frac{1}{2}\,g^{\lambda\rho}
  \!\left[
    \frac{\partial g_{\rho\nu}}{\partial\sigt}\nabla_\mu\sigt
  + \frac{\partial g_{\rho\mu}}{\partial\sigt}\nabla_\nu\sigt
  - \frac{\partial g_{\mu\nu}}{\partial\sigt}\nabla_\rho\sigt
  \right].
  }
  \label{eq:QGD_connection}
\end{equation}
This is the complete Levi-Civita connection, now expressed as a linear functional of
$\nabla_\mu\sigt$ scaled by the metric's structural sensitivity $\partial g/\partial\sigt$.
What GR calls spacetime curvature is, in QGD, the kinematic gradient flow of the
$\sigma$-field.

\subsection{The Advection Tensor}

For the two-body problem, the total field is $\sigt^{\rm tot} = \sigt^{(1)} + \sigt^{(2)}$,
so $\nabla_\mu\sigt^{\rm tot} = \nabla_\mu\sigt^{(1)} + \nabla_\mu\sigt^{(2)}$.
Substituting into~\eqref{eq:QGD_connection}, the connection of the effective one-body
system splits as:
\begin{equation}
  \widetilde\Gamma^\lambda_{\mu\nu}
  = \mathring\Gamma^\lambda_{\mu\nu}\!\left[\sigt^{(2)}\right]
  + \Omega^\lambda_{\mu\nu}\!\left[\sigt^{(1)},\sigt^{(2)}\right],
  \label{eq:connection_split}
\end{equation}
where $\mathring\Gamma^\lambda_{\mu\nu}[\sigt^{(2)}]$ is the background connection
generated by body~2 alone (driving the standard Keplerian orbit), and the
\emph{Advection Tensor} is:

\begin{definition}[Advection Tensor]
\label{def:advection}
\begin{equation}
  \boxed{
  \Omega^\lambda_{\mu\nu}
  = \frac{1}{2}\,g^{\lambda\rho}
  \!\left[
    \frac{\partial g_{\rho\nu}}{\partial\sigt}\,\nabla_\mu\sigt^{(1)}
  + \frac{\partial g_{\rho\mu}}{\partial\sigt}\,\nabla_\nu\sigt^{(1)}
  - \frac{\partial g_{\mu\nu}}{\partial\sigt}\,\nabla_\rho\sigt^{(1)}
  \right].
  }
  \label{eq:advection}
\end{equation}
\end{definition}

\begin{remark}
The Advection Tensor is \emph{exact}---it follows from the chain rule
applied to $g_{\mu\nu}(\sigt^{(1)} + \sigt^{(2)})$ and does not require any
perturbative approximation.  The factor $\nu$ appearing in the post-Newtonian force
law is not intrinsic to $\Omega^\lambda_{\mu\nu}$ itself; it enters via the
Tamir--Mino energy mapping when the two-body system is reduced to an effective
one-body description (Chapter~5).  The notes' representation of $\Omega$ as
``$\propto \nu$'' is an approximation valid only after the EOB reduction.
\end{remark}

\subsection{Explicit components in the Schwarzschild sector}

For the composite Schwarzschild metric $g_{tt} = -(1-\sigt^2)$, $g_{rr} = 1$
(isotropic gauge), the only non-vanishing metric sensitivity is:
\begin{equation}
  \frac{\partial g_{tt}}{\partial\sigt} = 2\sigt,
  \qquad
  \frac{\partial g_{ij}}{\partial\sigt} = 0.
  \label{eq:metric_sensitivity}
\end{equation}
The inverse metric components are $g^{tt} = -1/(1-\sigt^2)$ and $g^{rr} = 1$.
From Definition~\ref{def:advection}, the non-vanishing Advection Tensor components
for a static test body~1 ($\nabla_t\sigt^{(1)} = 0$) are:
\begin{align}
  \Omega^t_{tr}
  &= \frac{1}{2}\,g^{tt}\cdot 2\sigt^{(2)}\cdot\nabla_r\sigt^{(1)}
  = -\frac{\sigt^{(2)}\,\nabla_r\sigt^{(1)}}{1-(\sigt^{(2)})^2},
  \label{eq:Omega_ttr}\\
  \Omega^r_{tt}
  &= -\frac{1}{2}\,g^{rr}\cdot 2\sigt^{(2)}\cdot\nabla_r\sigt^{(1)}
  = -\sigt^{(2)}\,\nabla_r\sigt^{(1)}.
  \label{eq:Omega_rtt}
\end{align}
Substituting the single-body fields $\sigt^{(2)} = \sqrt{r_s^{(2)}/r}$ and
$\nabla_r\sigt^{(1)} = -\tfrac{1}{2}\sqrt{r_s^{(1)}}\,r^{-3/2}$:
\begin{equation}
  \Omega^r_{tt}
  = \frac{1}{2}\,\sqrt{\frac{r_s^{(1)}\,r_s^{(2)}}{r^3}}
  = \frac{1}{2}\,\frac{\sqrt{4G^2 m_1 m_2/c^4}}{r^{3/2}}
  \propto \frac{\sqrt{m_1 m_2}}{r^2}.
  \label{eq:Omega_rtt_explicit}
\end{equation}
After the Tamir--Mino mapping $\sqrt{m_1 m_2} \to \mu\sqrt{\nu}\,M$, this becomes
a $\nu$-weighted correction to the Newtonian force, generating the Ladder~A PN terms.

\subsection{The modified geodesic equation: explicit form}

Substituting~\eqref{eq:connection_split} into the geodesic condition, the equation
of motion for body~1 in the field of body~2 is:
\begin{equation}
  \frac{d^2 x^\mu}{ds^2}
  + \mathring\Gamma^\mu_{\alpha\beta}\,\frac{dx^\alpha}{ds}\frac{dx^\beta}{ds}
  = -\Omega^\mu_{\alpha\beta}\,\frac{dx^\alpha}{ds}\frac{dx^\beta}{ds}.
  \label{eq:modified_geodesic_explicit}
\end{equation}
The left side is the standard geodesic equation in the field of body~2; the right
side is the Advection force---the resistance of body~1's $\sigma$-field to being
``dragged'' through body~2's gradient.

In the static, non-relativistic limit with $dx^r/ds \ll 1$ and $dx^t/ds \approx 1$:
\begin{equation}
  \ddot r \approx
  -\frac{GM_2}{r^2}
  \underbrace{- \Omega^r_{tt}}_{\text{Advection force}}
  = -\frac{GM_2}{r^2}
  + \sigt^{(2)}\,\nabla_r\sigt^{(1)}.
  \label{eq:NR_limit}
\end{equation}
At Newtonian order, the second term is $\mathcal{O}(r_s^2/r^3)$ — a higher-order
PN correction — confirming that the advection force is invisible at Solar-System
scales and becomes significant only in the strong-field regime.

%=============================================================
\section{Stress-Energy Conservation}
\label{sec:conservation}
%=============================================================

\subsection{The stress-energy tensors}

In QGD three distinct stress-energy contributions arise from the master action:

\paragraph{Matter stress-energy.}
This is the standard $T^{\mu\nu}_{\rm matter}$, derived from varying
$\mathcal{L}_{\rm eff}(\psi, g(\sigma))$ with respect to the metric.

\paragraph{$\sigma$-field stress-energy.}
From the kinetic term in the action, contraction of the master field
equation~\eqref{eq:PU_field} with $\sigma^\mu$ via the Noether procedure gives:
\begin{equation}
  T^{\alpha\beta}_{\rm field}
  = \nabla^\alpha\sigma_\mu\,\nabla^\beta\sigma^\mu
  - \frac{1}{2}\,g^{\alpha\beta}\,(\nabla\sigma)^2
  + \kappa\lQ^2\text{-corrections},
  \label{eq:T_field}
\end{equation}
which is positive-definite in the exterior region and locally conserved.

\paragraph{Pais--Uhlenbeck stress-energy.}
The massive mode $G_{\mQ}$ carries its own stress-energy through the fourth-order
kinetic sector:
\begin{equation}
  \Theta^{\mu\nu}_{\rm PU}
  \sim \kappa\lQ^2
  \!\left[
    \nabla^2\sigma_\rho\,\nabla^\mu\nabla^\nu\sigma^\rho
    - \frac{1}{2}\,g^{\mu\nu}\,\nabla^2\sigma_\rho\,\nabla^2\sigma^\rho
  \right].
  \label{eq:T_PU}
\end{equation}
This tensor carries the repulsive pressure associated with the massive graviton mode
and is responsible for the UV regulation of the theory near $r \sim \lQ$.

\subsection{The full conservation law}

The total stress-energy conservation is:
\begin{equation}
  \boxed{
  \nabla_\nu\!\left(
    T^{\mu\nu}_{\rm matter}
    + T^{\mu\nu}_{\rm field}
    + \Theta^{\mu\nu}_{\rm PU}
  \right) = 0.
  }
  \label{eq:total_conservation}
\end{equation}
In the test-particle limit ($T^{\mu\nu}_{\rm matter} = \mu\,u^\mu u^\nu\,\delta^{(3)}$),
equation~\eqref{eq:total_conservation} reduces to the modified geodesic
equation~\eqref{eq:modified_geodesic_explicit}, with the divergence of
$T^{\mu\nu}_{\rm field}$ supplying the Advection force on the right side.

In the quantum-gravity regime ($r \sim \lQ$), the $\Theta^{\mu\nu}_{\rm PU}$ term
contributes a pressure that prevents the total stress-energy from diverging at a
point: the conservation law~\eqref{eq:total_conservation} holds everywhere, including
at $r = 0$, with $\Theta^{\mu\nu}_{\rm PU}$ providing the necessary compensating
pressure.

\subsection{Reduction to standard GR}

\begin{proposition}[GR limit]
\label{prop:GR_limit}
In the limits $\lQ \to 0$ and $|\sigma| \ll 1$ (weak field):
\begin{enumerate}[nosep]
  \item $\Theta^{\mu\nu}_{\rm PU} \to 0$, so matter and field stress-energies
    are conserved separately.
  \item $\Omega^\lambda_{\mu\nu} \to 0$, recovering the standard geodesic equation.
  \item $T^{\mu\nu}_{\rm field} = T^{\mu\nu}_{\rm GW}$ (standard gravitational
    wave stress-energy tensor).
  \item Equation~\eqref{eq:total_conservation} reduces to $\nabla_\nu T^{\mu\nu} = 0$
    (Einstein's conservation law).
\end{enumerate}
\end{proposition}

The corrections are suppressed by $\kappa\lQ^2/r^2 \sim 10^{-70}\,\mathrm{m}^{-2}/r^2$
at astrophysical scales, consistent with all current tests of GR.

%=============================================================
\section{Comparison: GR versus QGD Equations of Motion}
\label{sec:comparison}
%=============================================================

\begin{center}
\renewcommand{\arraystretch}{1.5}
\begin{tabular}{p{4.5cm}p{5cm}p{5cm}}
\toprule
\textbf{Feature} & \textbf{General Relativity} & \textbf{QGD} \\
\midrule
Fundamental law &
  Geodesic equation &
  $\sigma$-field equation \\
Connection &
  Levi-Civita $\Gamma[\partial g]$ &
  Composite $\Gamma[(\partial g/\partial\sigma)\nabla\sigma]$ \\
Advection term &
  Absent &
  $\Omega^\lambda_{\mu\nu}[\sigma_1, \sigma_2]$ \\
Particle nature &
  Test mass on background &
  Source of the $\sigma$-field \\
Stress conservation &
  $\nabla_\nu T^{\mu\nu}_{\rm m} = 0$ &
  $\nabla_\nu(T_{\rm m} + T_{\rm field} + \Theta_{\rm PU})^{\mu\nu} = 0$ \\
High-curvature limit &
  Singular at $r = 0$ &
  Regulated by $\Theta^{\mu\nu}_{\rm PU}$ \\
Force law (radial) &
  $-GM/r^2 \times [1 + \text{PN}]$ &
  $-GM/r^2 \times [1 - 3\nu u^2 + \tfrac{63707}{720}\nu u^3]$ \\
Weak-field limit &
  Newton's law &
  Newton's law (identical) \\
\bottomrule
\end{tabular}
\end{center}

%=============================================================
\section{The QGD Force Law}
\label{sec:force_law}
%=============================================================

\subsection{From the connection to the radial force}

The radial component of the modified geodesic equation~\eqref{eq:modified_geodesic_explicit},
projected along the separation vector, gives the effective radial force law.  In the
post-Newtonian expansion, using $u \equiv GM/rc^2$ as the small parameter:

\begin{theorem}[QGD radial force law]
\label{thm:force_law}
The effective radial acceleration in the QGD two-body system, derived from the
Advection-corrected geodesic equation and expressed in terms of the EOB potential
$A(u)$ (Chapter~5), is:
\begin{equation}
  \boxed{
  a_r = -\frac{GM}{r^2}
  \left[
    1
    \underbrace{-\,3\nu u^2\vphantom{\frac{X}{X}}}_{\text{Ladder A}}
    \underbrace{+\,\frac{63707}{720}\,\nu u^3}_{\text{Ladder B}}
  \right]
  + \mathcal{O}(u^4).
  }
  \label{eq:force_law}
\end{equation}
\end{theorem}

\begin{proof}
The EOB potential $A(u) = 1 - 2u + 2\nu u^3 - (63707/1440)\nu u^4$ is derived
from the QGD integral solution in Chapter~5.  The effective radial force is
$a_r = -\tfrac{1}{2}(GM/r^2)(dA/du)$, which gives:
\begin{equation}
  \frac{dA}{du} = -2 + 6\nu u^2 - \frac{63707}{360}\nu u^3.
\end{equation}
Multiplying by $-\tfrac{1}{2}$ yields~\eqref{eq:force_law}.  \qquad
\end{proof}

\subsection{Physical interpretation of each term}

\paragraph{Newtonian term ($1$).}
At $u \to 0$ (large separations), the brackets reduce to unity and the force
recovers Newton's inverse-square law exactly.  This is mandatory: both the
Levi-Civita connection and the Advection Tensor vanish in the Newtonian limit.

\paragraph{Ladder A term ($-3\nu u^2$, attractive).}
This is the leading PN correction.  The coefficient $-3\nu$ arises from the
cross-source $Q_t^{\rm cross} = 2\nabla\sigma_1\cdot\nabla\sigma_2$, which is the
$\sigma$-field analogue of the graviton self-interaction.  It deepens the potential
well, driving perihelion precession and the post-Newtonian inspiral.  In the GR
limit ($G_0$-only propagator), this is the standard 1PN force.

\paragraph{Ladder B term ($+63707\nu u^3/720$, repulsive).}
This term arises exclusively from the massive propagator $G_{\mQ}$ in the QGD
dual-propagator $G_0 - G_{\mQ}$.  Its positive sign means it acts as a repulsive
``quantum pressure'' that opposes the Ladder A attraction at high curvatures.  At
$u \sim 0.15$ (near the ISCO), the Ladder B contribution is approximately 30\%
of the Ladder A term for $\nu = 1/4$.

\subsection{Hierarchy of scales}

The three terms in~\eqref{eq:force_law} become comparable at different curvatures:

\begin{itemize}[nosep]
  \item Newtonian dominates for $u \ll 1$ (all astrophysical scales).
  \item Ladder A is $\mathcal{O}(\epsilon^2)$ where $\epsilon = v/c = \sqrt{u}$:
    significant for tight binaries and LIGO/LISA sources.
  \item Ladder B is $\mathcal{O}(\epsilon^3)$: a 1.5PN correction relative to
    Ladder A, detectable in principle with next-generation detectors.
\end{itemize}

The force law~\eqref{eq:force_law} is thus the connecting expression between the
abstract $\sigma$-field dynamics (Chapter~3) and the observable gravitational wave
phasing (Chapter~5).

%=============================================================
\section{Summary}
\label{sec:summary}
%=============================================================

\begin{center}
\renewcommand{\arraystretch}{1.4}
\begin{tabular}{p{6.5cm}p{6.5cm}}
\toprule
Equation & Meaning \\
\midrule
$\boxg(1-\kappa\lQ^2\boxg)\sigma_\mu = S_\mu$ &
  Field-level EoM: all motion encoded here \\
$g_{\mu\nu} = g_{\mu\nu}(\sigma_t)$ &
  Metric compositeness: no geometry without field \\
$\partial_\mu g_{\alpha\beta} = (\partial g_{\alpha\beta}/\partial\sigma_t)\,\nabla_\mu\sigma_t$ &
  Chain rule: curvature = $\sigma$-gradient flow \\
$\Gamma^\lambda_{\mu\nu} = \tfrac{1}{2}g^{\lambda\rho}[(\partial g/\partial\sigma)\nabla\sigma]^{\rm sym}$ &
  QGD composite connection \\
$\Omega^\lambda_{\mu\nu} = \Gamma^\lambda[\sigma_{\rm tot}] - \Gamma^\lambda[\sigma_2]$ &
  Advection Tensor: field-dragging correction \\
$\ddot x^\mu + (\Gamma^\mu + \Omega^\mu)_{\alpha\beta}\dot x^\alpha\dot x^\beta = 0$ &
  Modified worldline equation \\
$\nabla_\nu(T_{\rm m} + T_{\rm field} + \Theta_{\rm PU})^{\mu\nu} = 0$ &
  Total stress-energy conservation \\
$a_r = -(GM/r^2)[1 - 3\nu u^2 + (63707/720)\nu u^3]$ &
  QGD force law (to 4PN) \\
\bottomrule
\end{tabular}
\end{center}

\paragraph{What is established.}
\begin{itemize}[nosep]
  \item The QGD connection is derived exactly from the particle action via the
    composite-metric chain rule~\eqref{eq:metric_chain}.
  \item The Advection Tensor $\Omega^\lambda_{\mu\nu}$ is defined precisely
    (Definition~\ref{def:advection}) and its components computed explicitly
    (equations~\eqref{eq:Omega_ttr}--\eqref{eq:Omega_rtt_explicit}).
  \item The $\nu$-dependence in the force law does not originate from $\Omega$
    itself but from the Tamir--Mino EOB energy mapping.
  \item The three stress-energy contributions and their conservation law are
    identified (equation~\eqref{eq:total_conservation}).
  \item The GR limit is recovered exactly (Proposition~\ref{prop:GR_limit}).
\end{itemize}

\end{document}
